%% annex/p04/degrees-and-priority.tex
%% Parallel sketch annex for migrated v1.1 relative-computation results and
%% the registered Tier-1 finite-injury method.

\begin{annexitem}{annex:p4.degrees.1}{thm:p04-relativized-adequacy}{Relativized adequacy}
\sketchnote{omitted: transition-table and simulation invariants inherited from
the source universal and parameter sketches.}
Relativize enumeration: role-preserving renaming, surrogate starts, the
binary-alphabet argument, and designated-slot normalization include query
rows.  Universal simulation forwards queries.  Parameterization copies query
rows without consulting \(B\), yielding one total computable
\(\mathsf{par}_n^m\) uniformly for every oracle.
\end{annexitem}

\begin{annexitem}{annex:p4.degrees.2}{thm:p04-jump}{Jump theorem}
\enlargethispage{4\baselineskip}
\sketchnote{the occurrence and universal-simulation clauses of relativized
adequacy remain sketches; the diagonal calculation itself is retained.}
The jump \(B'\) is the domain of \(e\mapsto\mathsf{Univ}^B(e,e)\), so it is c.e.
relative to \(B\).  Suppose an oracle-\(B\) machine decided \(B'\), with
total characteristic function \(h^B\).  Define the oracle-computable partial
function
\[
 d^B(e)\simeq
 \begin{cases}
  0,&h^B(e)=0,\\
  \text{undefined},&h^B(e)=1.
 \end{cases}
\]
Relativized adequacy supplies \(b\) with \(\varphi^B_b\simeq d^B\).  Then
\[
 b\in B'
 \Leftrightarrow \varphi^B_b(b)\downarrow
 \Leftrightarrow h^B(b)=0
 \Leftrightarrow b\notin B',
\]
a contradiction.  Hence \(B'\not\leq_TB\).
\end{annexitem}

\begin{annexitem}{annex:p4.degrees.3}{thm:p04-jump-convention-bridge}{Convention bridge for the jump}
\sketchnote{omitted: transition-table proofs for four transformations; the
monotonicity reduction still depends on the sketched relativized parameter
compiler and padding.}
Query-free embedding and negative answers for oracle \(\varnothing\) give
\(K\sim_m\varnothing'\).  For halting versus definedness, respectively compile
``on any halt write \(0\) and accept'' and ``loop unless acceptance has valid
numerical output''; hence the jumps are many-one equivalent.

If \(A\leq_TB\), replace \(A\)-queries by oracle-\(B\) decisions; relativized
parameterization and padding give \(A'\leq_mB'\).  Finally compile, from \(n\),
a program accepting with valid output \(0\) exactly when \(n\in B\) and
otherwise diverging.  Its total index map reduces \(B\) to \(B'\).  The jump
theorem gives strictness, so Post-theorem statements survive the convention
switch.
\end{annexitem}

\begin{annexitem}{annex:p4.degrees.4}{thm:p04-friedberg-muchnik}{Friedberg--Muchnik theorem}
\sketchnote{omitted: a primitive-recursive coding of finite stage states and
the fair scheduler; the formal proof that the effective stage approximations
and enumeration of oracle functionals agree with the live numbering; the
coding of an empty restraint for a computation making no query; the uniform
index calculation for the two c.e. stage unions; and the compiler identifying
every alleged reduction with a requirement in the displayed list.  The
strategy actions, restraint convention, priority induction, and satisfaction
verification are retained.}
List
\[
 R_e:\varphi_e^B\not\simeq\operatorname{char}_A,
 \qquad S_e:\varphi_e^A\not\simeq\operatorname{char}_B
\]
in the priority order \(R_0>S_0>R_1>S_1>\cdots\).  Use the effective stage
approximations \(A_s,B_s\) and the oracle-program numbering fixed immediately
before Definition~\ref{def:p04-turing-jump}.  Write
\(\varphi_{e,s}^{B_s}(x)\mathord\downarrow=y\) when the computation with the
finite oracle approximation \(B_s\) halts within \(s\) steps with output
\(y\).  A fair computable schedule visits
each requirement in priority order at arbitrarily late stages.  An initialized
\(R_e\) chooses a number \(x\) larger than all active higher-priority
restraints and distinct from every reserved witness, and protects
\(x\notin A\).  It waits for a stage computation
\(\varphi_{e,s}^{B_s}(x)\downarrow=0\).  Its query use \(u\) is the largest
oracle position queried; if no query occurs, the restraint is empty.  On
seeing such a computation, it puts \(x\) into \(A\), preserves every
\(B\)-oracle answer below \(u+1\), and initializes every lower-priority strategy.
Initialization discards a witness and restraint but never removes a number
already enumerated.  A lower-priority strategy is forbidden to enumerate
inside a protected oracle segment unless a higher-priority action first
injures it.  If the final stabilized computation returns \(1\), keeping
\(x\notin A\) already gives disagreement; if it never converges, disagreement
holds by partiality.  The \(S_e\)-strategy interchanges \(A\) and \(B\).

\enlargethispage{5\baselineskip}
Induct down the priority list.  Once all higher-priority strategies have
reached their final outcomes, they enumerate only finitely many further
numbers and impose fixed restraints.  The current strategy is therefore
initialized only finitely often.  Its final witness is fresh above those
restraints.  If it acts, its final oracle use is preserved and its enumeration
flips the computed value \(0\); if it observes a stable value \(1\), the
protected witness stays out; if it waits forever, the computation diverges.
Thus every \(R_e\) and \(S_e\) is met, giving
\(A\not\leq_TB\) and \(B\not\leq_TA\).  The omitted numbering compiler is
the bookkeeping step ensuring that every candidate oracle reduction occurs
among these requirements.

The stage instructions effectively enumerate \(A\) and \(B\), so both are
c.e.; the omitted uniform index calculation formalizes this observation.
Neither
is computable, since every computable set is reducible to every oracle.  Every
c.e. set is reducible to \(\varnothing'\); if either \(A\) or \(B\) had degree
\(\mathbf0'\), the other would reduce to it, contradicting incomparability.
Hence both degrees lie strictly between \(\mathbf0\) and \(\mathbf0'\).
\end{annexitem}

\begin{annexitem}{annex:p4.degrees.5}{thm:p04-many-one-versus-turing}{Many-one versus Turing reducibility}
\enlargethispage{4\baselineskip}
\begin{proof}
If \(f\) witnesses \(A\leq_mB\), an oracle-\(B\) machine computes \(f(x)\)
and makes one membership query, so \(A\leq_TB\).  For the converse,
\(\N\setminus K\leq_TK\) by querying \(K\) and reversing the answer.  Were
\(\N\setminus K\leq_mK\), the computable-enumerability clause of
Theorem~\ref{thm:p04-reduction-preservation} and c.e.-ness of \(K\) would make
\(\N\setminus K\) c.e., contrary to
Theorem~\ref{thm:p04-k-complement}.
\end{proof}
\end{annexitem}
